%====================================================================
%  Directional Congestion and Tidal Lanes
%
%  A self-contained statement of the directional quantitative spatial
%  commuting model: setup, equilibrium definition, master inventory of
%  variables/parameters/data sources, and the estimation pipeline
%  (lambda from a directed speed-flow regression; inversion of
%  fundamentals; nested-fixed-point exact-hat solver for counterfactuals).
%====================================================================
\documentclass[12pt]{article}

\usepackage[margin=1in]{geometry}
\usepackage{amsmath,amssymb,amsthm,mathtools}
\usepackage{graphicx,booktabs,tabularx,array,setspace,float,enumitem}
\usepackage{longtable}
\usepackage{natbib}
\usepackage[colorlinks=true,linkcolor=blue,citecolor=blue,urlcolor=blue]{hyperref}
\setstretch{1.15}
\setlength{\parindent}{1.5em}

\newtheorem{proposition}{Proposition}
\newtheorem{lemma}{Lemma}
\newtheorem{definition}{Definition}
\newtheorem{assumption}{Assumption}
\newtheorem{remark}{Remark}

\DeclareMathOperator*{\argmax}{arg\,max}

\newcommand{\N}{\mathcal N}
\newcommand{\E}{\mathcal E}
\newcommand{\R}{\mathcal R}
\newcommand{\Lbar}{\bar L}
\newcommand{\Wbar}{\bar W}
\newcommand{\ubar}{\bar u}
\newcommand{\Abar}{\bar A}
\newcommand{\tbar}{\bar t}
\newcommand{\tkl}{t_{kl}}
\newcommand{\tijk}{\tau_{ij}}
\newcommand{\obs}{\mathrm{obs}}
\newcommand{\cf}{\mathrm{cf}}

\title{Directional Congestion and Tidal Lanes:\\
Model Part}
\author{Yige Duan \and Xiaoruo Feng \and Yizhen Gu}
\date{\today}

\begin{document}
\maketitle

\begin{abstract}
\noindent
This paper develops a quantitative spatial commuting model in which every primitive of the road network is directed. Lane capacities, free-flow speeds, link travel times, and bilateral commuting cost indices are allowed to differ between $(k,l)$ and $(l,k)$; workers choose a residence, a workplace, and a route under Fréchet preferences; and equilibrium delivers gravity in commuting flows, gravity in directed link traffic, and a closed-form welfare scale. Three features distinguish this paper. The first is directional speed asymmetry at the segment level: observed peak-period imbalances between the two directions of a corridor enter the cost structure directly, rather than being subsumed in a symmetric travel-time index. The second is the policy experiment: tidal-lane reallocation reassigns lane capacity between directions within existing road width, requiring neither new land supply nor substantial physical construction, making it among the lowest-cost levers available to urban congestion management. The third is the data input: granular, high-frequency segment-level speed observations identify the direction-specific congestion wedge on each corridor and supply the directed cost primitives the model requires.
\end{abstract}

\tableofcontents
\clearpage

%====================================================================
\section{Setup}\label{sec:setup}
%====================================================================

\subsection{Locations and directed network}

The city is partitioned into $N$ spatial cells $\N=\{1,\dots,N\}$,
each indexed by its centroid.  The road network is a directed graph
$(\N,\E)$ with edge set $\E\subseteq\N\times\N$.  A directed edge
$(k,l)\in\E$ records an admissible move from cell $k$ to cell $l$.
$(k,l)\in\E$ does not imply $(l,k)\in\E$, and even when both edges
exist the two directions are distinct objects with distinct
primitives.

\subsection{Directed link primitives}

Each directed link $(k,l)\in\E$ carries three primitives:
\begin{itemize}[leftmargin=2em,itemsep=0.2em]
  \item physical length $\ell_{kl}>0$ (centroid-to-centroid);
  \item directed lane capacity $n_{kl}>0$;
  \item free-flow speed $v^{0}_{kl}>0$, with associated free-flow
        travel time $\tbar_{kl}\equiv\ell_{kl}/v^{0}_{kl}$.
\end{itemize}
Realised (congested) directed travel time $\tkl$ is endogenous and
determined by the congestion technology in
\S\ref{sec:objects}.  Directionality is permitted in every primitive:
$\ell_{kl}\ne\ell_{lk}$, $n_{kl}\ne n_{lk}$, $v^{0}_{kl}\ne v^{0}_{lk}$,
$\tkl\ne t_{lk}$ are all admitted.

\subsection{Iceberg cost convention}\label{ssec:iceberg}

The model uses a dimensionless iceberg link cost $\tkl\ge 1$ with the
\citet{AllenArkolakis2022} micro-foundation
\begin{equation}
  \tkl\;=\;\bigl(\ell_{kl}/v_{kl}\bigr)^{1/\theta},
  \qquad
  v_{kl}\;=\;v^{0}_{kl}\bigl(n_{kl}/\xi_{kl}\bigr)^{\lambda},
  \label{eq:iceberg}
\end{equation}
where $v_{kl}$ is the realised (congested) directed speed and
$\xi_{kl}$ is endogenous directed link traffic.  The iceberg cost
$\tkl$ enters utility multiplicatively as
$\prod_m t_{r_{m-1},r_m}$ and the kernel of the route system is
$\tkl^{-\theta}$.  Working with $\tkl$ rather than physical minutes
guarantees scale invariance of the equilibrium under a change of time
units; physical minutes enter only through the empirical free-flow
and congested speeds, $v^{0}_{kl}$ and $v_{kl}$.

\subsection{Population and fundamentals}

A unit mass $\Lbar$ of workers is allocated across residences
$\{L_i^R\}$ and workplaces $\{L_j^F\}$ with
$\sum_i L_i^R=\sum_j L_j^F=\Lbar$.  Every cell is endowed with a
baseline residential amenity $\ubar_i>0$ and a baseline workplace
productivity $\Abar_j>0$.  Endogenous amenity and productivity load
the local scale of activity through
\begin{equation}
  u_i\;=\;\ubar_i\bigl(L_i^R\bigr)^{\beta},
  \qquad
  A_j\;=\;\Abar_j\bigl(L_j^F\bigr)^{\alpha},
  \qquad \alpha>0,\ \beta<0.
  \label{eq:fundamentals}
\end{equation}
Firms operate under perfect competition with labour as the only
input, so the market-clearing wage at workplace $j$ is
\begin{equation}
  w_j\;=\;A_j\;=\;\Abar_j\bigl(L_j^F\bigr)^{\alpha}.
  \label{eq:wage}
\end{equation}
The composite amenity $u_i$ absorbs the equilibrium effect of
housing-market clearing; goods trade and non-traded amenities are
not modelled separately, consistent with the Allen--Arkolakis
reduced-form parameterisation.

\subsection{Parameters}

The model is closed by four scalar parameters:
\begin{itemize}[leftmargin=2em,itemsep=0.2em]
  \item $\theta>0$, the Fr\'echet shape parameter governing route and
        residence--workplace dispersion;
  \item $\alpha>0$, the production externality (agglomeration)
        elasticity in \eqref{eq:fundamentals};
  \item $\beta<0$, the residential externality (congestion)
        elasticity in \eqref{eq:fundamentals};
  \item $\lambda\ge 0$, the congestion elasticity of speed with
        respect to traffic per lane in \eqref{eq:iceberg}.
\end{itemize}
$(\theta,\alpha,\beta)$ are externally calibrated; $\lambda$ is
internally estimated from a directed speed--flow regression
(\S\ref{ssec:lambda_est}).

%====================================================================
\section{Equilibrium objects}\label{sec:objects}
%====================================================================

This section assembles the equations used in the equilibrium
definition.  Workers choose a residence, a workplace, and a route;
firms hire under perfect competition; the network technology maps
directed link traffic into directed link costs.  Each object is
introduced in turn.

\paragraph{Indirect utility.}
A worker $\nu$ chooses a residence $i$, a workplace $j$, and a route
$r=(r_0=i,r_1,\dots,r_K=j)$ with $(r_{m-1},r_m)\in\E$:
\begin{equation}
  V_{ij,r}(\nu)
  \;=\;
  \frac{u_i\,w_j}{\prod_{m=1}^{K(r)}t_{r_{m-1},r_m}}\,
  \varepsilon_{ij,r}(\nu),
  \qquad
  \varepsilon_{ij,r}\overset{\mathrm{iid}}{\sim}\mathrm{Fr\acute echet}(\theta).
  \label{eq:utility}
\end{equation}

\paragraph{Bilateral commuting cost index.}
Sum over feasible routes:
\begin{equation}
  \tijk^{-\theta}
  \;\equiv\;
  \sum_{r\in\R_{ij}}\prod_{m=1}^{K(r)}t_{r_{m-1},r_m}^{-\theta}.
  \label{eq:tau_def}
\end{equation}
Collect link costs in the directed weighted adjacency matrix
$A=[a_{kl}]$ with $a_{kl}=\tkl^{-\theta}\mathbf 1\{(k,l)\in\E\}$.
Under the spectral-radius assumption below, the series converges and
\begin{equation}
  \tijk^{-\theta}\;=\;\bigl[(I-A)^{-1}\bigr]_{ij}.
  \label{eq:tau_series}
\end{equation}
The bilateral cost index is generically asymmetric,
$\tijk\ne\tau_{ji}$.

\begin{assumption}[Spectral radius]\label{ass:specrad}
$\rho(A)<1$, where $\rho(\cdot)$ is the spectral radius.
\end{assumption}

\paragraph{Aggregate welfare scale.}
Define
\begin{equation}
  \Wbar\;\equiv\;
  \Bigl[\sum_{(i,j)\in\N^2}\tijk^{-\theta}\,u_i^{\theta}\,w_j^{\theta}\Bigr]^{1/\theta}.
  \label{eq:Wbar_def}
\end{equation}
Ex-ante expected welfare is $\mathbb W=\Gamma(1-1/\theta)\Wbar$.

\paragraph{Gravity equation.}
The bilateral commuting flow is
\begin{equation}
  L_{ij}\;=\;\Lbar\cdot
  \frac{\tijk^{-\theta}\,u_i^{\theta}\,w_j^{\theta}}{\Wbar^{\theta}},
  \qquad
  L_i^R=\sum_j L_{ij},\quad L_j^F=\sum_i L_{ij}.
  \label{eq:gravity}
\end{equation}

\paragraph{Market-access indices.}
Define
\begin{equation}
  \Pi_i^{-\theta}\;\equiv\;\sum_{j}\tijk^{-\theta}\,L_j^F\,P_j^{\theta},
  \qquad
  P_j^{-\theta}\;\equiv\;\sum_{i}\tijk^{-\theta}\,L_i^R\,\Pi_i^{\theta}.
  \label{eq:MA_def}
\end{equation}
Substituting into \eqref{eq:gravity} delivers the equivalent
market-access form of gravity,
\begin{equation}
  L_{ij}\;=\;\tijk^{-\theta}\cdot
  \frac{L_i^R}{\Pi_i^{-\theta}}\cdot
  \frac{L_j^F}{P_j^{-\theta}}.
  \label{eq:gravity_MA}
\end{equation}
Because $\tijk$ is directional, $\Pi_i\ne P_i$ in general.

\paragraph{Link-traffic gravity.}
Directed link traffic $\xi_{kl}$ is the expected number of
traversals of edge $(k,l)$ summed over commuters:
\begin{equation}
  \xi_{kl}\;=\;\tkl^{-\theta}\,P_k^{-\theta}\,\Pi_l^{-\theta}.
  \label{eq:xi_gravity}
\end{equation}
Each component is directional, so $\xi_{kl}\ne\xi_{lk}$ in general.

\paragraph{Congestion technology.}
The iceberg cost in \eqref{eq:iceberg}, after substitution of
$v_{kl}=v^{0}_{kl}(n_{kl}/\xi_{kl})^{\lambda}$, takes the compact
form
\begin{equation}
  \tkl
  \;=\;
  \tbar_{kl}^{1/\theta}\,
  \bigl(\xi_{kl}/n_{kl}\bigr)^{\lambda/\theta}.
  \label{eq:cong}
\end{equation}
Substituting \eqref{eq:xi_gravity} into \eqref{eq:cong} and solving
for $\tkl$ delivers the reduced-form expression
\begin{equation}
  \tkl
  \;=\;
  \tbar_{kl}^{1/[\theta(1+\lambda)]}\cdot
  n_{kl}^{-\lambda/[\theta(1+\lambda)]}\cdot
  \bigl(P_k\Pi_l\bigr)^{-\lambda/(1+\lambda)}.
  \label{eq:tkl_reduced}
\end{equation}

\paragraph{Tidal-lane channel.}
A tidal-lane policy reallocates lanes on a target set
$\E^{\star}\subseteq\E$ from off-peak to peak,
\begin{equation}
  n_{kl}^{\cf}=n_{kl}+\Delta_{kl},\qquad
  n_{lk}^{\cf}=n_{lk}-\Delta_{kl},\qquad
  (k,l)\in\E^{\star},
\end{equation}
preserving total corridor capacity $n_{kl}+n_{lk}$.  Because each
direction enters \eqref{eq:cong} separately, $\tkl$ falls and
$t_{lk}$ rises strictly on treated corridors; the induced changes in
$\{\tijk,L_{ij},L^R,L^F\}$ are asymmetric across $(i,j)$.  Within an
undirected model the policy is vacuous.

\paragraph{Population shares.}
Let $l_i^R\equiv L_i^R/\Lbar$ and $l_j^F\equiv L_j^F/\Lbar$.
Combining \eqref{eq:fundamentals}, \eqref{eq:wage},
\eqref{eq:gravity}, and the marginal-population identities yields
\begin{align}
  (l_i^R)^{1-\beta\theta}
  &\;=\;\chi\sum_{j}\tijk^{-\theta}\,\ubar_i^{\theta}\,\Abar_j^{\theta}\,
       (l_j^F)^{\alpha\theta},
  \label{eq:eq_R}\\
  (l_j^F)^{1-\alpha\theta}
  &\;=\;\chi\sum_{i}\tijk^{-\theta}\,\ubar_i^{\theta}\,\Abar_j^{\theta}\,
       (l_i^R)^{\beta\theta},
  \label{eq:eq_F}
\end{align}
with the scalar $\chi\equiv\Lbar^{\theta(\alpha+\beta)}/\Wbar^{\theta}$
pinned down by the normalisation
$\prod_i\ubar_i=\prod_j\Abar_j=1$.

%====================================================================
\section{Definition of Equilibrium}\label{sec:eq}
%====================================================================

\begin{definition}[Equilibrium]\label{def:eq}
Given primitives $\{(\ubar_i,\Abar_j)\}_{i,j\in\N}$ and
$\{(\ell_{kl},v^{0}_{kl},n_{kl})\}_{(k,l)\in\E}$, parameters
$(\theta,\alpha,\beta,\lambda)$, and total population $\Lbar$, an
equilibrium is the tuple
\[
  \mathcal Q
  \;\equiv\;
  \bigl(
    \{L_i^R\},\,\{L_j^F\},\,\{L_{ij}\},\,
    \{\tkl\},\,\{\tijk\},\,\{\xi_{kl}\},\,
    \{\Pi_i\},\,\{P_j\},\,\Wbar
  \bigr)
\]
that satisfies, jointly,
\begin{enumerate}[label=(\textbf{E\arabic*}),leftmargin=3em,itemsep=0.25em]
\item \textbf{Bilateral cost index.} \eqref{eq:tau_series}:
  $\tijk^{-\theta}=[(I-A)^{-1}]_{ij}$, with
  $A_{kl}=\tkl^{-\theta}\mathbf 1\{(k,l)\in\E\}$.
\item \textbf{Gravity flows.} \eqref{eq:gravity}:
  $L_{ij}=\Lbar\,\tijk^{-\theta}u_i^{\theta}w_j^{\theta}/\Wbar^{\theta}$,
  with $u_i,w_j$ from \eqref{eq:fundamentals}--\eqref{eq:wage} and
  $\Wbar$ from \eqref{eq:Wbar_def}.
\item \textbf{Market clearing.} $L_i^R=\sum_j L_{ij}$,
  $L_j^F=\sum_i L_{ij}$, $\sum_i L_i^R=\sum_j L_j^F=\Lbar$.
\item \textbf{Market-access indices.} \eqref{eq:MA_def} for
  $(\Pi_i,P_j)$.
\item \textbf{Link-traffic gravity.} \eqref{eq:xi_gravity}:
  $\xi_{kl}=\tkl^{-\theta}P_k^{-\theta}\Pi_l^{-\theta}$.
\item \textbf{Congestion technology.} \eqref{eq:cong}:
  $\tkl=\tbar_{kl}^{1/\theta}(\xi_{kl}/n_{kl})^{\lambda/\theta}$.
\item \textbf{Spectral-radius condition.} Assumption~\ref{ass:specrad}.
\end{enumerate}
Equivalently, equilibrium share allocations $(l^R,l^F)$ solve the
fixed-point system
\eqref{eq:eq_R}--\eqref{eq:eq_F} jointly with
\eqref{eq:tkl_reduced}.
\end{definition}

\begin{remark}[Variable count and equation count]
Conditions (E1)--(E6) deliver
$|\N|^2$ equations for $\{\tijk\}$,
$|\N|^2$ equations for $\{L_{ij}\}$,
$2|\N|$ for $\{L_i^R,L_j^F\}$,
$2|\N|$ for $\{\Pi_i,P_j\}$,
$|\E|$ for $\{\xi_{kl}\}$, and
$|\E|$ for $\{\tkl\}$, matching the dimension of the unknowns;
$\Wbar$ is fixed by either of the welfare identities
$\Wbar^{\theta}=\sum_i u_i^{\theta}\Pi_i^{-\theta}=\sum_j
w_j^{\theta}P_j^{-\theta}$ from \eqref{eq:Wbar_def}.
\end{remark}

\paragraph{Existence and uniqueness.}
Existence under Assumption~\ref{ass:specrad} and $\lambda\ge 0$
follows from the Brouwer-based fixed-point argument in Proposition~1
of \citet{AllenArkolakis2022}.  Their Proposition~1(b) supplies a
sufficient uniqueness condition,
\begin{equation}
  \alpha\;\le\;\tfrac{1}{2}\bigl(1/\theta-\lambda\bigr),
  \qquad
  \beta\;\le\;\tfrac{1}{2}\bigl(1/\theta-\lambda\bigr).
  \label{eq:AAunique}
\end{equation}
Under our sign restrictions $\alpha>0,\beta<0$ the right-hand
inequality is implied, and the left-hand inequality is binding.  At
the baseline calibration $(\theta,\alpha)=(6.83,0.06)$,
\eqref{eq:AAunique} requires $\lambda\le 1/\theta-2\alpha\approx
0.026$, a value below the central estimate of $\lambda$ (see
\S\ref{ssec:lambda_est}).  Where \eqref{eq:AAunique} fails, no
formal claim of uniqueness is made.  Three numerical diagnostics are
reported instead on every grid point in the counterfactual output:
multi-start invariance of \textsc{HatSolve}, the Jacobian spectral
radius at the converged solution, and smoothness of the equilibrium
in the underlying parameters.

\paragraph{Welfare representation.}
Two equivalent forms hold at any equilibrium:
\begin{equation}
  \Wbar^{\theta}
  \;=\;\sum_i u_i^{\theta}\Pi_i^{-\theta}
  \;=\;\sum_j w_j^{\theta}P_j^{-\theta}.
  \label{eq:welfare_MA}
\end{equation}
Equation \eqref{eq:welfare_MA} is the only welfare object the model
identifies; counterfactual welfare changes are constructed from
exact-hat algebra (\S\ref{ssec:hat}) and reported as the time-equivalent
variation $\hat{\mathbb W}$.  No separate housing-market or goods-market
welfare component is computed.

%====================================================================
\section{Variables, parameters, and data sources}\label{sec:inventory}
%====================================================================

This section is the backbone of the implementation: every model
object is classified by (i) role (data primitive, parameter,
inverted fundamental, endogenous equilibrium variable), (ii) source
(direct construction from data, external calibration, internal
estimation, equilibrium inversion, equilibrium computation), and
(iii) whether it is held fixed in counterfactual analysis (a
\emph{state} variable in the policy sense).

\subsection{Classification convention}\label{ssec:classification}

\begin{itemize}[leftmargin=2em,itemsep=0.2em]
\item \textbf{Data primitive (P)}: constructed directly from raw
  data; held fixed in baseline and varies only by assumption in
  counterfactuals.  Examples: $\ell_{kl}$, $n_{kl}$.
\item \textbf{External parameter (X)}: a scalar borrowed from the
  literature with documented external estimates; not re-estimated on
  our data in the baseline.  Examples: $\theta$, $\alpha$, $\beta$.
\item \textbf{Internally estimated parameter (I)}: a scalar
  identified from a moment in the data.  Example: $\lambda$.
\item \textbf{Inverted fundamental (V)}: a vector recovered from
  the equilibrium system to rationalise the observed allocation
  given the data and parameters.  Examples: $\ubar_i,\Abar_j$.
\item \textbf{Endogenous equilibrium variable (Q)}: a vector
  determined by Definition~\ref{def:eq} given primitives,
  parameters, and inverted fundamentals.  Examples: $L_{ij},\tijk,
  \xi_{kl},\tkl,\Pi_i,P_j,\Wbar$.
\item \textbf{Counterfactual policy lever ($\Delta$)}: a primitive
  perturbed by assumption in a counterfactual.  Example: $n_{kl}^{\cf}$.
\end{itemize}
The state space of the policy problem is the data primitives plus
the inverted fundamentals: $(\ell,n,v^0,\ubar,\Abar)$ and the
parameters $(\theta,\alpha,\beta,\lambda)$.  Endogenous equilibrium
variables are recomputed at every counterfactual policy.

\subsection{Master inventory}\label{ssec:inventory}

Table~\ref{tab:inventory} lists every model object that appears in
Definition~\ref{def:eq}, along with its empirical realisation and
its data or estimation source.

\begin{small}
\begin{longtable}{p{2.0cm}p{4.4cm}cp{6.3cm}}
\caption{Model variables, parameters, and data sources.}
\label{tab:inventory}\\
\toprule
Symbol & Object & Type & Source / construction \\
\midrule
\endfirsthead
\multicolumn{4}{l}{\emph{(continued)}}\\
\toprule
Symbol & Object & Type & Source / construction \\
\midrule
\endhead
\midrule
\multicolumn{4}{r}{\emph{(continued on next page)}}\\
\endfoot
\bottomrule
\endlastfoot
\multicolumn{4}{l}{\textbf{A. Spatial primitives}}\\
$i,j\in\N$ & Spatial cell (node) & P & Square baseline (3{,}357 cells, mean area 8.38\,km$^2$); H3 hex (4{,}604 cells, 6.10\,km$^2$) and network-adapted Voronoi (1{,}213 cells, 23.20\,km$^2$) as robustness.\\
$(k,l)\in\E$ & Directed grid edge & P & Adjacency in the largest weakly connected component supported by the grid travel-cost map of \S\ref{ssec:rawgrid_step}.\\
$\ell_{kl}$ & Directed grid length & P & Centroid-to-access-point distance plus on-network road-path length; pair-specific shortest path on the directed road graph (edge-projected version, \S\ref{ssec:rawgrid_step}).\\
$n_{kl}$ & Directed lane count & P / $\Delta$ & Length-weighted lane proxy from a roadtype-to-lanes lookup on matched centerlines (mean 3.13, median 3.32, length coverage 92.0\%); perturbed in counterfactuals.\\
$v^{0}_{kl}$ & Free-flow directed speed & P & Time-weighted harmonic mean of matched-centerline speeds in 22:00--05:00 (citywide mean 38.46\,km/h).\\
$v^{\obs}_{kl}$ & Observed congested speed & P & Same aggregator at the AM peak 07:00--09:00 (mean 35.19\,km/h); PM peak 17:00--19:00 (33.33\,km/h) used as the second period.\\
$\tbar_{kl}$ & Free-flow time & P & $\ell_{kl}/v^{0}_{kl}$.\\
$\tkl^{\obs}$ & Observed iceberg link cost & P & $(\ell_{kl}/v_{kl}^{\obs})^{1/\theta}$ with $v_{kl}^{\obs}$ length-weight projected from the matched centerlines onto the grid edge.\\
\midrule
\multicolumn{4}{l}{\textbf{B. Population primitives}}\\
$L_{ij}^{\obs}$ & Bilateral commuting flow & P & 100\,m grid OD matrix point-in-polygon aggregated to the active tessellation; Beijing 2022 baseline (6.91M raw OD pairs over 453{,}391 unique 100\,m grids); Shenzhen 2024 (3.97M pairs over 110{,}339 grids) for replication.\\
$L_i^{R,\obs}$ & Residence marginal & P & Row sum $\sum_j L_{ij}^{\obs}$.\\
$L_j^{F,\obs}$ & Workplace marginal & P & Column sum $\sum_i L_{ij}^{\obs}$.\\
$\Lbar$ & Total commuters & P & $\sum_i L_i^{R,\obs}=\sum_j L_j^{F,\obs}$ on the largest weakly connected component of $(\N,\E)$.\\
\midrule
\multicolumn{4}{l}{\textbf{C. Parameters}}\\
$\theta$ & Fr\'echet shape & X & Calibrated to $6.83$ from \citet{Ahlfeldt2015}; gravity grid-search of \S\ref{ssec:theta_check} reported as a robustness diagnostic.\\
$\alpha$ & Production externality elasticity & X & Calibrated to $0.06$, central estimate of \citet{CombesGobillon2015}; grid $\{0.04,0.06,0.08\}$.\\
$\beta$ & Residential externality elasticity & X & Calibrated to $-0.30$ following \citet{AllenArkolakis2014,MonteReddingRossiHansberg2018,HeblichReddingSturm2020}; grid $\{-0.50,-0.30,-0.20\}$.\\
$\lambda$ & Congestion elasticity & I & Internally estimated by IV speed--flow regression \eqref{eq:lambda_regression}; counterfactual robustness over $\lambda\in\{0.1,0.3,0.6,1.0\}$.\\
\midrule
\multicolumn{4}{l}{\textbf{D. Inverted fundamentals}}\\
$\ubar_i$ & Baseline residential amenity & V & Recovered from \eqref{eq:invert_R}--\eqref{eq:invert_F} given $(\theta,\alpha,\beta)$, $(L_i^{R,\obs},L_j^{F,\obs})$, and $\tijk^{\obs}$.\\
$\Abar_j$ & Baseline workplace productivity & V & Recovered jointly with $\ubar_i$, normalised so that $\prod_i\ubar_i=\prod_j\Abar_j=1$.\\
\midrule
\multicolumn{4}{l}{\textbf{E. Endogenous equilibrium variables}}\\
$L_i^R, L_j^F$ & Residence / workplace populations & Q & Solve \eqref{eq:eq_R}--\eqref{eq:eq_F}; equal to $L_i^{R,\obs},L_j^{F,\obs}$ at baseline by Assumption~\ref{ass:obs_eq}.\\
$L_{ij}$ & Bilateral commuting flow & Q & \eqref{eq:gravity}; equals $L_{ij}^{\obs}$ at baseline by Assumption~\ref{ass:obs_eq}.\\
$u_i,w_j$ & Composite amenity / wage & Q & \eqref{eq:fundamentals}--\eqref{eq:wage}.\\
$\tijk$ & Bilateral cost index & Q & \eqref{eq:tau_series}; baseline $\tijk^{\obs}$ uses $A^{\obs}_{kl}=(\tkl^{\obs})^{-\theta}$ with $\tkl^{\obs}$ converted to iceberg form via \eqref{eq:iceberg}.\\
$\xi_{kl}$ & Directed link traffic & Q & \eqref{eq:xi_gravity} at the equilibrium $(P_k,\Pi_l,\tkl)$; baseline empirical analogue $\hat\xi_{kl}^{\obs}$ in \S\ref{ssec:link_traffic_obs}.\\
$\Pi_i, P_j$ & Resident / firm market access & Q & Linear system \eqref{eq:MA_def}.\\
$\tkl$ & Directed iceberg link cost & Q & \eqref{eq:cong}; baseline $\tkl^{\obs}$ converted from $v_{kl}^{\obs}$ via the iceberg micro-foundation \eqref{eq:iceberg}.\\
$\Wbar$ & Aggregate welfare scale & Q & \eqref{eq:Wbar_def} or \eqref{eq:welfare_MA}; counterfactual $\hat{\Wbar}$ in \eqref{eq:hat_W}.\\
\end{longtable}
\end{small}

\subsection{Where each parameter is identified}\label{ssec:identification_map}

\begin{itemize}[leftmargin=2em,itemsep=0.2em]
\item $\theta$ enters the model in three places: (i) the kernel
  $\tkl^{-\theta}$ of the cost-index series, (ii) the iceberg
  conversion exponent in \eqref{eq:iceberg}, and (iii) the
  externality bands $\theta\alpha,\theta\beta$.  External calibration
  is the baseline because the gravity-based internal estimator
  \eqref{eq:gravity_log} requires inversion of $\tijk(\theta)$ at
  every grid point and is reported as a robustness check rather than
  as the identifying moment.
\item $\alpha$ and $\beta$ enter only the share equations
  \eqref{eq:eq_R}--\eqref{eq:eq_F} through the externality powers
  $L^{F,\alpha\theta}$ and $L^{R,\beta\theta}$.  In a single
  cross-section without an exogenous shifter of $L^R$ or $L^F$ they
  are not separately identified from $(\ubar_i,\Abar_j)$; external
  calibration is the only credible route.
\item $\lambda$ has a single identifying moment in the data, the
  speed--flow regression \eqref{eq:lambda_regression}.  This is the
  one parameter the paper estimates internally.
\item $(\ubar_i,\Abar_j)$ are exactly identified by the marginal
  populations $(L_i^{R,\obs},L_j^{F,\obs})$ and the bilateral cost
  index $\tijk^{\obs}$ under Assumption~\ref{ass:obs_eq}.  This is
  the inversion step.
\end{itemize}

\subsection{Welfare measurement: reduced form, not full estimation}
\label{ssec:welfare_strategy}

The model identifies a single welfare scalar $\Wbar$ via
\eqref{eq:welfare_MA}, equivalently
$\mathbb W=\Gamma(1-1/\theta)\Wbar$.  Welfare changes in
counterfactuals are reported as $\hat{\mathbb W}=\hat{\Wbar}$,
computed in closed form from the exact-hat formula
\eqref{eq:hat_W}.  Three consequences follow.
\begin{enumerate}[label=(\roman*),leftmargin=2em,itemsep=0.2em]
\item No housing-market or goods-market block is estimated
  separately; the composite amenity $u_i$ absorbs the equilibrium
  effect of housing-market clearing under the AA-class
  parameterisation.
\item Welfare changes are interpreted in time-equivalent variation:
  a counterfactual with $\hat{\mathbb W}=1.02$ is equivalent in
  welfare terms to a uniform 1.96\% reduction in every bilateral
  commuting time.
\item No estimate of the marginal value of time, of housing supply
  elasticity, or of the elasticity of substitution across goods is
  required; the welfare aggregator is fully pinned down by
  $(\theta,\alpha,\beta,\lambda)$ and the baseline shares.
\end{enumerate}

%====================================================================
\section{Estimation and Algorithms}\label{sec:est}
%====================================================================

The empirical pipeline executes in five steps:
\begin{enumerate}[label=\textbf{Step \arabic*.},leftmargin=3.5em,itemsep=0.2em]
\item Construct the directed primitives
  $(\ell_{kl},v^{0}_{kl},v_{kl}^{\obs},n_{kl})$ and the population
  primitives $(L_{ij}^{\obs},L_i^{R,\obs},L_j^{F,\obs},\Lbar)$
  (\S\ref{ssec:data_construction}).
\item Calibrate $(\theta,\alpha,\beta)$ externally
  (\S\ref{ssec:calibration}).
\item Estimate $\lambda$ from the IV speed--flow regression
  (\S\ref{ssec:lambda_est}).
\item Invert $(\ubar_i,\Abar_j)$ from the equilibrium system at the
  observed allocation (\S\ref{ssec:invert}).
\item Solve counterfactuals via the nested-fixed-point exact-hat
  algorithm and read off $\hat{\mathbb W}$ (\S\ref{ssec:hat}).
\end{enumerate}
Each step is described below as inputs $\to$ algorithm $\to$ outputs.

\subsection{Step 1: Constructing data inputs}\label{ssec:data_construction}

The data primitives in Table~\ref{tab:inventory} are built from two
raw inputs through a six-stage pipeline.  This subsection states the
construction for the directed grid network and the population
matrix.

\paragraph{Raw inputs.}
The first input is a 100\,m grid origin--destination commuting
matrix.  Beijing 2022 contains 6{,}909{,}608 OD pairs over
453{,}391 unique 100\,m grids; Shenzhen 2024 contains 3{,}970{,}569
OD pairs over 110{,}339 grids.  The matrix records counts of
home--work commuters between grid pairs.  The second input is a
directed segment-level speed panel observed at 10-minute frequency
on the same cities, with raw fields \texttt{cityname},
\texttt{roadseg\_id}, \texttt{roadname}, \texttt{roadtype},
\texttt{semantic}, \texttt{location}, and \texttt{geometry}.  In the
retained Beijing run the panel contributes 23.29M raw speed
observations after quality control.

\paragraph{Raw-to-centerline matching.}
The raw road map is collapsed to a single \emph{raw-centerline}
system in two parts.  A manual corridor channel handles 963 raw
edges (3.9\% of all raw edges, 20.4\% of total raw length) on long
divided expressways, frontage roads, and grade-separated
interchanges where mechanical skeletonisation is unreliable; the
remaining edges enter a skeleton branch built with a 50\,m road
buffer, a 5\,m raster, and 50\,m short-branch pruning.  Each
undirected skeleton centerline is duplicated into two directions.
Every raw sub-segment is then matched to a directed centerline
candidate within a 60/180/600\,m search radius using
\[
  \mathrm{score}(r,c)\;=\;d_{\mathrm{mean}}(r,c)+0.1\,\Delta\theta(r,c),
\]
and accepted only if $d_{\mathrm{mean}}(r,c)\le 180\,\text{m}$.
After matching, 99.42\% of baseline-kept split segments and 95.59\%
of all raw edges are matched, accounting for 98.73\% of total raw
length.  The retained baseline directed network has 11{,}890
directed centerlines spanning 19{,}391\,km.

\paragraph{Directed centerline speeds.}
For each baseline directed centerline $c$ and time bin $\tau$
(weekday status by hour), the speed aggregator is the time-weighted
harmonic mean
\begin{equation}
  \bar v_{c,\tau}^{\,\mathrm{time}}
  \;=\;
  \frac{\sum_{r\in\R(c,\tau)}\ell_r}
       {\sum_{r\in\R(c,\tau)}\ell_r/v_r},
  \label{eq:harmonic_speed}
\end{equation}
with $\R(c,\tau)$ the matched raw segments for $c$ in $\tau$.  The
retained run yields 394{,}468 directed centerline-hour cells across
8{,}349 unique directed centerlines (citywide weekday means: AM
peak 35.19\,km/h, PM peak 33.33\,km/h, 22:00--05:00 free-flow
window 38.46\,km/h).  The morning peak 07:00--09:00 is the baseline
period; the evening peak 17:00--19:00 is retained for the
second-period checks in CF1--CF4.

\paragraph{Lane-count proxy.}
A roadtype-to-lanes lookup is applied to the matched raw segments
and length-weighted to the directed centerline.  Coverage on the
baseline directed network is 70.6\% by count and 92.0\% by length,
with mean lane estimate 3.13, median 3.32, and a mean
opposite-direction lane ratio of 1.04 (86.2\% of corridors symmetric
within $\pm 0.5$).

\paragraph{Spatial tessellation $\N$.}
Three grid systems discretise the study area.  The square baseline
contains 3{,}357 cells with mean area 8.38\,km$^2$.  An H3 hexagonal
grid (4{,}604 cells, 6.10\,km$^2$) and a network-adapted Voronoi
tessellation (1{,}213 cells, 23.20\,km$^2$, generated from
baseline-centerline node seeds under a 1.5\,km minimum-distance rule
and a 200\,km$^2$ cell-area cap) are retained as robustness systems.
The OD matrix is point-in-polygon assigned from 100\,m to the active
tessellation,
$L_{ij}^{\obs}=\sum_{(g,h)\in i\times j}L^{100\mathrm{m}}_{gh}$, and
the reachability filter keeps every $(i,j)$ whose endpoints both lie
in the largest weakly connected component of $(\N,\E)$.

\paragraph{Grid-to-grid travel cost.}\label{ssec:rawgrid_step}
The bilateral cost map between adjacent grid pairs is computed on
the directed road graph rather than on the centerline backbone.
Among three candidate constructions evaluated on AM adjacent square
pairs --- the centerline benchmark, a node-attached version that
snaps each grid centroid to the nearest road node, and an
edge-projected version that projects each centroid onto a nearby
directed road edge --- the operating construction is the
edge-projected version: each grid centroid is projected onto a
nearby directed road edge, virtual HOME and WORK nodes connect only
to the active home and work cells, and the pair-specific shortest
path through the restricted graph defines $\ell_{kl}$ and the
on-network minute count.  Edge projection raises connected-pair
coverage to 0.474 (square) / 0.494 (hex) / 0.875 (Voronoi) and
keeps connector distances at 831\,m / 850\,m / 432\,m on the three
grids.  The centerline benchmark is retained as a second metric,
used in the asymmetry screen of \S\ref{ssec:scenarios} and in
robustness diagnostics.

\paragraph{Aggregating to directed grid links.}
For each baseline directed centerline piece $p$ falling on grid edge
$(k,l)$, the pipeline records the in-grid length $\ell_p$, the
AM-peak harmonic-mean speed $v_p^{\obs}$, the 22:00--05:00 free-flow
speed $v^{0}_{p}$, and the lane-count proxy $n_p$.  Aggregating to
the directed grid link gives
\begin{align}
  \ell_{kl}
    &\;=\;\sum_{p\in P(k,l)}\ell_{p},
  \label{eq:agg_ell}\\
  v_{kl}^{\obs}
    &\;=\;\Bigl(\sum_{p\in P(k,l)}\ell_p\Bigr)\Big/
       \Bigl(\sum_{p\in P(k,l)}\ell_p / v_p^{\obs}\Bigr),
  \label{eq:agg_v}\\
  v^{0}_{kl}
    &\;=\;\Bigl(\sum_{p\in P(k,l)}\ell_p\Bigr)\Big/
       \Bigl(\sum_{p\in P(k,l)}\ell_p / v^{0}_p\Bigr),\\
  n_{kl}
    &\;=\;\Bigl(\sum_{p\in P(k,l)}\ell_p\,n_p\Bigr)\Big/
       \Bigl(\sum_{p\in P(k,l)}\ell_p\Bigr),
\end{align}
where $P(k,l)$ is the set of centerline pieces on grid edge $(k,l)$.
The aggregator \eqref{eq:agg_v} matches \eqref{eq:harmonic_speed} at
the centerline level.

\paragraph{Iceberg link costs.}
The iceberg link costs at baseline follow the convention of
\S\ref{ssec:iceberg}:
\begin{equation}
  \tkl^{\obs}\;=\;(\ell_{kl}/v_{kl}^{\obs})^{1/\theta},
  \qquad
  \tbar_{kl}^{\,\mathrm{ice}}\;=\;(\ell_{kl}/v^{0}_{kl})^{1/\theta}.
  \label{eq:tkl_obs_iceberg}
\end{equation}
The baseline cost index $\tijk^{\obs}$ follows from
$A^{\obs}_{kl}=(\tkl^{\obs})^{-\theta}$ via \eqref{eq:tau_series}.
Pieces missing a lane proxy or an AM speed observation, which are
concentrated on short peripheral fragments rather than on the
arterial backbone, inherit the centerline-mean lane count and the
roadtype-conditional median speed; the affected directed grid links
are flagged and dropped in robustness checks.

\paragraph{Population primitives.}
With $L_{ij}^{\obs}$ from the Stage 5 point-in-polygon assignment,
\[
  L_i^{R,\obs}=\sum_j L_{ij}^{\obs},
  \qquad
  L_j^{F,\obs}=\sum_i L_{ij}^{\obs},
  \qquad
  \Lbar=\sum_i L_i^{R,\obs}.
\]

\paragraph{Constructed link traffic $\hat\xi_{kl}^{\obs}$.}\label{ssec:link_traffic_obs}
Directed flows are not directly observed.  Empirical link traffic is
constructed by route assignment of $L_{ij}^{\obs}$ over $(\N,\E)$
using the model-implied route probabilities
\begin{equation}
  \pi^{kl}_{ij}\;=\;\frac{\tau_{ik}^{-\theta}\,\tkl^{-\theta}\,\tau_{lj}^{-\theta}}
                          {\tijk^{-\theta}},
  \qquad
  \hat\xi_{kl}^{\obs}\;=\;\sum_{i,j}L_{ij}^{\obs}\,\pi^{kl}_{ij},
  \label{eq:hat_xi_obs}
\end{equation}
evaluated with the iceberg link costs $\tkl^{\obs}$ from
\eqref{eq:tkl_obs_iceberg}.  If the observed allocation were an
equilibrium, $\hat\xi_{kl}^{\obs}$ would coincide with
\eqref{eq:xi_gravity}.

\subsection{Step 2: External calibration of $(\theta,\alpha,\beta)$}
\label{ssec:calibration}

\begin{itemize}[leftmargin=2em,itemsep=0.2em]
\item $\theta=6.83$ \citep{Ahlfeldt2015}.
\item $\alpha=0.06$ \citep{CombesGobillon2015,CombesDurantonGobillon2008,AllenArkolakis2022}.
\item $\beta=-0.30$ \citep{AllenArkolakis2014,MonteReddingRossiHansberg2018,HeblichReddingSturm2020}.
\end{itemize}
Robustness grids $\alpha\in\{0.04,0.06,0.08\}$ and
$\beta\in\{-0.50,-0.30,-0.20\}$ are reported alongside the main
estimates.  An internal check on $\theta$ via the gravity grid
search of \S\ref{ssec:theta_check} is reported as a diagnostic.

\subsection{Step 3: Estimating $\lambda$ from a directed
speed--flow regression}\label{ssec:lambda_est}

\paragraph{Estimating equation.}
Substituting the iceberg micro-foundation \eqref{eq:iceberg} into
its log form gives the identity
$\ln v_{kl}=\ln v^{0}_{kl}-\lambda\ln(\xi_{kl}/n_{kl})$.  Relaxing
the intercept to absorb link-type variation yields
\begin{equation}
  \boxed{\;
  \ln v_{kl}^{\obs}
  \;=\;
  m_0
  -\lambda\,\ln\!\bigl(\hat\xi_{kl}/n_{kl}\bigr)
  +\eta_{\mathrm{bin}(kl)}
  +\varepsilon_{kl}.
  \;}
  \label{eq:lambda_regression}
\end{equation}
$\eta_{\mathrm{bin}(kl)}$ absorbs systematic capacity differences
across lane bins and roadtype classes; standard errors are clustered
at the centerline parent.

\paragraph{Constructing the regressor.}
Under the naive assignment, $\hat\xi_{kl}$ uses $\tkl^{\obs}$ which
is itself a function of the dependent variable $v_{kl}^{\obs}$.  To
break the mechanical circularity, two non-IV variants are reported,
together with the IV estimator that is the baseline:
\begin{itemize}[leftmargin=2em,itemsep=0.15em]
\item \emph{Free-flow regressor (a)}: replace $\{\tkl^{\obs}\}$ by
  $\{\tbar_{kl}^{\,\mathrm{ice}}\}$ in $\pi^{kl}_{ij}$, so that
  $\hat\xi_{kl}^{\mathrm{ff}}$ depends only on nighttime free-flow
  costs.
\item \emph{Leave-one-out regressor (b)}: recompute $\pi^{kl}_{ij}$
  with link $(k,l)$ excluded from the cost index.
\item \emph{Naive regressor (c)}: $\hat\xi_{kl}$ from
  $\{\tkl^{\obs}\}$, reported only as a diagnostic of the
  circularity bias.
\end{itemize}

\paragraph{IV strategy.}
The baseline estimator instruments $\ln(\hat\xi_{kl}/n_{kl})$ with
link-geography primitives that shift directed traffic per lane
without affecting contemporaneous speed shocks.  Candidate
instruments are
(i) historical lane counts predating the current commuting pattern,
(ii) link-to-CBD distance,
(iii) the reverse-direction nighttime free-flow speed on the same
corridor.  The first stage regresses $\ln(\hat\xi_{kl}/n_{kl})$ on
$z_{kl}$ and the lane-bin fixed effects; the second stage uses the
fitted value in \eqref{eq:lambda_regression}.  Identification rests
on the exclusion restriction that $z_{kl}$ affects $\ln v_{kl}$ only
through traffic per lane, conditional on the fixed effects.

\paragraph{Output.}
The pipeline returns $\hat\lambda^{\mathrm{IV}}$ together with the
first-stage $F$, the over-identification statistic when more than
one instrument is used, and the OLS variants (a)--(c) as
comparison rows.  Counterfactual robustness is conducted at
$\lambda\in\{0.1,0.3,0.6,1.0\}$ to bracket both the OLS range and
the AA estimate converted to our iceberg convention.

\subsection{Step 4: Inversion of $(\ubar_i,\Abar_j)$}\label{ssec:invert}

\begin{assumption}[Observed allocation is an equilibrium]
\label{ass:obs_eq}
$(L_i^{R,\obs},L_j^{F,\obs})$ equals $(L_i^R,L_j^F)$ in
Definition~\ref{def:eq} under primitives
$(\ell_{kl},v_{kl}^{0},n_{kl})$, parameters
$(\theta,\alpha,\beta,\lambda)$, and some baseline fundamentals
$(\ubar_i,\Abar_j)$.  The bilateral cost index used in the inversion
is $\tijk^{\obs}=[(I-A^{\obs})^{-1}]_{ij}^{-1/\theta}$ with
$A^{\obs}_{kl}=(\tkl^{\obs})^{-\theta}$ from \eqref{eq:tkl_obs_iceberg}.
\end{assumption}

Under Assumption~\ref{ass:obs_eq}, the equilibrium system
\eqref{eq:eq_R}--\eqref{eq:eq_F} evaluated at the data is
\begin{align}
  (l_i^{R,\obs})^{1-\beta\theta}
  &\;=\;\chi\sum_{j}(\tijk^{\obs})^{-\theta}\,\ubar_i^{\theta}\,\Abar_j^{\theta}\,
       (l_j^{F,\obs})^{\alpha\theta},
  \label{eq:invert_R}\\
  (l_j^{F,\obs})^{1-\alpha\theta}
  &\;=\;\chi\sum_{i}(\tijk^{\obs})^{-\theta}\,\ubar_i^{\theta}\,\Abar_j^{\theta}\,
       (l_i^{R,\obs})^{\beta\theta}.
  \label{eq:invert_F}
\end{align}

\paragraph{Algorithm \textsc{Invert} (iterative fixed point).}

\begin{enumerate}[leftmargin=2em,itemsep=0.2em]
\item \textbf{Initialise.} $\ubar_i^{(0)}=1$, $\Abar_j^{(0)}=1$,
  $\chi^{(0)}=1$.
\item \textbf{Bilateral cost index.} Compute $\tijk^{\obs}$ once
  from $A^{\obs}_{kl}=(\tkl^{\obs})^{-\theta}$ via a single sparse
  $(I-A^{\obs})^{-1}$ solve.
\item \textbf{Update $\ubar$.} For each $i$,
\begin{equation}
  \ubar_i^{(k+1),\theta}
  \;=\;
  \frac{(l_i^{R,\obs})^{1-\beta\theta}}
       {\chi^{(k)}\sum_{j}(\tijk^{\obs})^{-\theta}\,
                    \Abar_j^{(k),\theta}\,(l_j^{F,\obs})^{\alpha\theta}}.
\end{equation}
\item \textbf{Update $\Abar$.} For each $j$,
\begin{equation}
  \Abar_j^{(k+1),\theta}
  \;=\;
  \frac{(l_j^{F,\obs})^{1-\alpha\theta}}
       {\chi^{(k)}\sum_{i}(\tijk^{\obs})^{-\theta}\,
                    \ubar_i^{(k+1),\theta}\,(l_i^{R,\obs})^{\beta\theta}}.
\end{equation}
\item \textbf{Normalise.} Rescale so that
  $\prod_i\ubar_i^{(k+1)}=\prod_j\Abar_j^{(k+1)}=1$;
  update $\chi^{(k+1)}$ from $\chi=\Lbar^{\theta(\alpha+\beta)}/
  \Wbar^{\theta}$.
\item \textbf{Convergence.} Stop when
  $\max_i|\ubar_i^{(k+1)}/\ubar_i^{(k)}-1|<10^{-8}$ and
  $\max_j|\Abar_j^{(k+1)}/\Abar_j^{(k)}-1|<10^{-8}$.
\end{enumerate}

\paragraph{Output.}
$(\ubar_i^{\star},\Abar_j^{\star})_{i,j\in\N}$ and the baseline
welfare scalar $\Wbar^{\star}$ from \eqref{eq:welfare_MA}; together
with the data primitives these fully pin down the baseline
equilibrium and are held fixed in counterfactual analysis.

\paragraph{Well-posedness.}
Under \eqref{eq:AAunique} and Assumption~\ref{ass:specrad}
\textsc{Invert} admits a unique strictly positive fixed point.
Outside this region, the three numerical diagnostics introduced
after Definition~\ref{def:eq} (multi-start invariance, Jacobian
spectral radius below one at the converged solution, and smoothness
in parameters) are reported on every grid point.

\subsection{Step 5: Counterfactual exact-hat solver}\label{ssec:hat}

For any endogenous $x$, write $\hat x\equiv x^{\cf}/x$.  The policy
shock is $\hat n_{kl}$; free-flow times do not change,
$\hat{\tbar}_{kl}=1$.  The exact-hat system follows by dividing each
equation in Definition~\ref{def:eq} by its baseline counterpart:
\begin{align}
  \hat\tkl^{\,1+\lambda}
  &\;=\;\hat n_{kl}^{-\lambda}\,(\hat P_k\hat\Pi_l)^{-\lambda},
  \label{eq:hat_tkl}\\
  \hat\tijk^{-\theta}
  &\;=\;
  \frac{[(I-A^{\cf})^{-1}]_{ij}}{[(I-A^{\obs})^{-1}]_{ij}},
  \quad A^{\cf}_{kl}=\tkl^{\obs,-\theta}\hat\tkl^{-\theta},
  \label{eq:hat_tau}\\
  \hat l_i^{R,1-\beta\theta}
  &\;=\;\hat\chi\cdot
    \frac{\sum_j (\tijk^{\obs})^{-\theta}\hat\tijk^{-\theta}
                  \Abar_j^{\star,\theta}(l_j^{F,\obs})^{\alpha\theta}
                  \hat l_j^{F,\alpha\theta}}
         {\sum_j (\tijk^{\obs})^{-\theta}\Abar_j^{\star,\theta}
                  (l_j^{F,\obs})^{\alpha\theta}},
  \label{eq:hat_lR}\\
  \hat l_j^{F,1-\alpha\theta}
  &\;=\;\hat\chi\cdot
    \frac{\sum_i (\tijk^{\obs})^{-\theta}\hat\tijk^{-\theta}
                  \ubar_i^{\star,\theta}(l_i^{R,\obs})^{\beta\theta}
                  \hat l_i^{R,\beta\theta}}
         {\sum_i (\tijk^{\obs})^{-\theta}\ubar_i^{\star,\theta}
                  (l_i^{R,\obs})^{\beta\theta}},
  \label{eq:hat_lF}\\
  \hat{\Wbar}^{\theta}
  &\;=\;\sum_{(i,j)}\pi_{ij}^{\obs}\,
        \hat l_i^{R,\beta\theta}\,\hat l_j^{F,\alpha\theta}\,
        \hat\tijk^{-\theta},
  \label{eq:hat_W}
\end{align}
with normalisations $\sum_i l_i^{R,\obs}\hat l_i^R=1$,
$\sum_j l_j^{F,\obs}\hat l_j^F=1$, $\hat\chi=\hat{\Wbar}^{-\theta}$,
and $\pi_{ij}^{\obs}=L_{ij}^{\obs}/\Lbar$.

\paragraph{Algorithm \textsc{HatSolve} (nested fixed point).}

\begin{enumerate}[leftmargin=2em,itemsep=0.2em]
\item \textbf{Inputs.} Baseline objects
  $\bigl(\tkl^{\obs},L_i^{R,\obs},L_j^{F,\obs},\tijk^{\obs},
  P_k,\Pi_l,\ubar_i^{\star},\Abar_j^{\star}\bigr)$ and counterfactual
  lane capacities $\{\hat n_{kl}\}$.
\item \textbf{Initialise.} $\hat\tkl=1$, $\hat l_i^R=1$,
  $\hat l_j^F=1$.
\item \textbf{Outer step on $\hat\tkl$.}
\begin{enumerate}[label=(\alph*),leftmargin=2em,itemsep=0.15em]
\item Form $A^{\cf}_{kl}=\tkl^{\obs,-\theta}\hat\tkl^{-\theta}$ and
  compute $\hat\tijk^{-\theta}$ from \eqref{eq:hat_tau} by a single
  sparse $(I-A^{\cf})^{-1}$ solve.
\item \textbf{Inner step on shares.} Iterate
  \eqref{eq:hat_lR}--\eqref{eq:hat_lF} jointly, with
  $\hat\chi=\hat{\Wbar}^{-\theta}$ from \eqref{eq:hat_W}, until
  $\max_i|\Delta\hat l_i^R|+\max_j|\Delta\hat l_j^F|<10^{-8}$.
\item Update market access: $\hat\Pi_i^{-\theta}$ and
  $\hat P_j^{-\theta}$ from the hat versions of \eqref{eq:MA_def}.
\item Update $\hat\tkl$ from \eqref{eq:hat_tkl}.
\end{enumerate}
\item \textbf{Convergence.} Stop the outer loop when
  $\max_{kl}|\hat\tkl^{(k+1)}/\hat\tkl^{(k)}-1|<10^{-6}$.
\item \textbf{Outputs.}
  $(\hat\tkl,\hat\xi_{kl},\hat\tijk,\hat l_i^R,\hat l_j^F,\hat{\Wbar})$
  and the welfare metric $\hat{\mathbb W}=\hat{\Wbar}$.
\end{enumerate}

\paragraph{Implementation notes.}
The bilateral-cost step is the binding cost; the sparse
$(I-A^{\cf})^{-1}$ solve is performed once per outer iteration.
Damped updates $\hat\tkl^{(k+1)}=\hat\tkl^{(k),\,1-\rho}\cdot
[\hat\tkl^{\mathrm{trial}}]^{\rho}$ with $\rho\in[0.3,0.7]$ stabilise
convergence at high $\lambda$; Anderson acceleration on the inner
share loop is reported as an option.  A levels
counterpart of \textsc{HatSolve} (re-running the equilibrium at
$(\ubar^{\star},\Abar^{\star})$ with counterfactual primitives) is
maintained as a numerical consistency check.

\subsection{Counterfactual scenarios}\label{ssec:scenarios}

\begin{description}[leftmargin=2.5em,itemsep=0.25em]
\item[CF1.] \emph{Symmetric vs.\ directional commuting costs.}
  Replace $\tkl^{\obs}$ by the corridor-level average
  $\tfrac12(\tkl^{\obs}+t_{lk}^{\obs})$, holding lane capacities and
  fundamentals fixed.  Isolates the welfare contribution of observed
  directionality.
\item[CF2.] \emph{Tidal-lane reallocation.} Apply
  $n_{kl}^{\cf}=n_{kl}+\Delta_{kl}$, $n_{lk}^{\cf}=n_{lk}-\Delta_{kl}$
  on a target set $\E^{\star}$ obtained from a directional-asymmetry
  screen on the centerline backbone.  For each centerline with valid
  peak-period speeds in both directions, define
  $\mathrm{ratio}_p=\min(v_{AB,p}/v_{BA,p},v_{BA,p}/v_{AB,p})$ and
  classify the centerline as \emph{asymmetric in peak $p$} if
  $\mathrm{ratio}_p<0.5$.  Three nested target sets are reported:
  the strict tidal-candidate set ($|\E^{\star}|=14$ corridors that
  also satisfy AM--PM reversal of the faster direction and
  $\mathrm{ratio}_p<0.5$ in both peaks), the AM-asymmetric set
  (405 corridors, 11.0\% of bidirectional centerlines), and the
  PM-asymmetric set (361 corridors, 9.8\%).  Lane shifts
  $\Delta_{kl}\in\{0.5,1,2\}$ preserve total corridor capacity
  $n_{kl}+n_{lk}$.
\item[CF3.] \emph{Uniform lane expansion.} Multiply $n_{kl}$ by a
  factor in $\{1.10,1.25,1.50\}$ uniformly across links.  Benchmarks
  CF2 against undifferentiated capacity addition.
\item[CF4.] \emph{Welfare and distributional comparisons.} Across
  CF1--CF3, report $\hat{\mathbb W}$, the cross-grid Gini of
  $(\hat l_i^R,\hat l_j^F)$, and the share of grids with
  $\hat l_i^R>1$.
\end{description}
The robustness array runs each scenario on the PM peak, the
hexagonal grid, $\theta\in\{4,6.83,9\}$, and
$\lambda\in\{0.1,0.3,0.6,1.0\}$.

\subsection{Internal robustness check on $\theta$}\label{ssec:theta_check}

For each $\theta\in\{3,4,\dots,12\}$ the procedure recomputes
$\tijk^{\obs}(\theta)$ via \eqref{eq:tau_series} and runs the
two-way fixed-effects gravity regression
\begin{equation}
  \ln L_{ij}^{\obs}\;=\;-\theta\,\ln\tijk(\theta)+\mu_i+\nu_j+\varepsilon_{ij}.
  \label{eq:gravity_log}
\end{equation}
The internal-consistency estimator selects the $\theta$ at which the
regression slope equals the imposed $\theta$.  The result is
reported alongside the external calibration; the external calibration
remains the baseline because \eqref{eq:gravity_log} relies on the
same network costs that are themselves equilibrium objects.

%====================================================================
\begin{thebibliography}{10}
%====================================================================
\bibitem[Ahlfeldt et al.(2015)]{Ahlfeldt2015}
Ahlfeldt, G., Redding, S. J., Sturm, D. M., Wolf, N. (2015). The
economics of density: Evidence from the Berlin Wall.
\emph{Econometrica}, 83(6), 2127--2189.

\bibitem[Allen and Arkolakis(2014)]{AllenArkolakis2014}
Allen, T., Arkolakis, C. (2014). Trade and the topography of the
spatial economy. \emph{Quarterly Journal of Economics}, 129(3),
1085--1140.

\bibitem[Allen and Arkolakis(2022)]{AllenArkolakis2022}
Allen, T., Arkolakis, C. (2022). The welfare effects of transportation
infrastructure improvements. \emph{Review of Economic Studies},
89(6), 2911--2957.

\bibitem[Combes, Duranton, and Gobillon(2008)]{CombesDurantonGobillon2008}
Combes, P.-P., Duranton, G., Gobillon, L. (2008). Spatial wage
disparities: Sorting matters! \emph{Journal of Urban Economics},
63(2), 723--742.

\bibitem[Combes and Gobillon(2015)]{CombesGobillon2015}
Combes, P.-P., Gobillon, L. (2015). The empirics of agglomeration
economies. In \emph{Handbook of Regional and Urban Economics}, Vol.
5A, ed.\ G.\ Duranton, J.\ V.\ Henderson, and W.\ Strange, 247--348.
Elsevier.

\bibitem[Duranton and Turner(2011)]{DurantonTurner2011}
Duranton, G., Turner, M. A. (2011). The fundamental law of road
congestion: Evidence from US cities. \emph{American Economic Review},
101(6), 2616--2652.

\bibitem[Heblich, Redding, and Sturm(2020)]{HeblichReddingSturm2020}
Heblich, S., Redding, S. J., Sturm, D. M. (2020). The making of the
modern metropolis: Evidence from London. \emph{Quarterly Journal of
Economics}, 135(4), 2059--2133.

\bibitem[Monte, Redding, and Rossi-Hansberg(2018)]{MonteReddingRossiHansberg2018}
Monte, F., Redding, S. J., Rossi-Hansberg, E. (2018). Commuting,
migration, and local employment elasticities. \emph{American Economic
Review}, 108(12), 3855--3890.

\end{thebibliography}

\end{document}
